\documentclass[10pt,a4paper]{amsart}
\usepackage[margin=19mm]{geometry}
\usepackage{amsmath,amssymb,mathtools}
\usepackage[scr=rsfs,cal=euler]{mathalfa}
\usepackage{xfrac}
\usepackage[unicode,hidelinks,bookmarks=false]{hyperref}
\newcommand{\ie}{i.e.\ }
\newcommand{\cf}{cf.\ }
\newcommand{\resp}{respectively\ }
\newcommand{\Subsection}{\subsection}
\providecommand{\nobreakdash}{\nobreak\hbox{-}}
\setlength{\emergencystretch}{2em}
\multlinegap0pt
\usepackage{bm}
\usepackage{bbm}
\usepackage{dsfont}
\usepackage{xspace}
\usepackage{cancel}
\usepackage{mathtools}
\usepackage{amsthm}
\makeatletter
\@ifundefined{lem}{\newtheorem{lem}{Lemma}}{}
\makeatother
\DeclareMathOperator{\spn}{span}
\newcommand{\An}{\mathfrak{A}}
\newcommand{\C}{\mathbb{C}}
\newcommand{\Oo}{\mathcal{O}}
\newcommand{\Restriction}{\upharpoonright}
\newcommand{\Sn}{\mathfrak{S}}
\newcommand{\Zz}{\mathcal{G}}
\newcommand{\Z}{\mathbb{Z}}
\newcommand{\rank}{\mathrm{rank}}
\newcommand{\rquotient}{\backslash}
\newcommand{\sets}[2]{\left\{#1\,\middle|\,#2\right\}}
\newcommand{\sgn}{\mathrm{sgn}}
\newcommand{\std}{\mathrm{std}}
\newcommand{\triv}{\mathrm{triv}}
\title[The asymptotic Hecke algebra and rigidity]{The asymptotic Hecke algebra and rigidity}
\author{Stefan Dawydiak}
\date{Source: arXiv:2312.11092v3 (CC BY 4.0)}
\begin{document}
\maketitle

\noindent\emph{Source and licence.} The asymptotic Hecke algebra and rigidity. Version arXiv:2312.11092v3 (CC BY 4.0), distributed under the Creative Commons Attribution 4.0 International licence, as recorded in the arXiv licence metadata for this exact version and shown on the abstract page https://arxiv.org/abs/2312.11092v3. Full rights notice: licenses/arxiv-2312.11092-CC-BY-4.0.txt.\par\smallskip

\noindent\textit{Editorial scope.} Counted unit: the Lem whose source label is \texttt{lem trho act with rank 1 or 0} in the article (source0.tex lines 886--891, source label \texttt{lem trho act with rank 1 or 0}), delivered together with the article's own complete original proof of it (source0.tex lines 893--1002, one \texttt{proof} environment of 110 lines). No mathematics is written, repaired or re-derived: statement and proof are the article's own text line for line. Required context retained uncounted: lem orthogonal idempotents for unextended set (lines 760--778); eqn trho mu transforms according to rho (lines 837--840). Every definition, display and label of those lines is retained unchanged. Omitted from this package and not redistributed: 1--759, 779--836, 841--885, 892--892, 1003--4976. Nothing inside a retained excerpt is omitted or altered: every retained body is delivered exactly as the article's lines are written, so no omission receipt of any kind accompanies this package. Citations: the retained excerpts carry no citation command at all, so no bibliography entry of the article is transcribed and no reference is redistributed. Reference adaptation: none was needed and none was made; every reference inside the delivered unit resolves inside it, so no unresolved reference or citation is produced. Printed slips kept literal: no printed word, symbol, hypothesis or number of the counted statement or of its proof was altered. Layout and class: the article's own class is \texttt{article} with packages including amsfonts, amsmath, amssymb, amsthm, amsrefs, arydshln, calc, graphicx, hyperref, geometry, todonotes, microtype, tikz-cd, dynkin-diagrams; the wrapper is the lane's pinned \texttt{amsart} editorial wrapper at 10pt on A4, which restates the theorem-like environments the retained text needs and adds the article's own macro definitions that the retained text needs (\texttt{\textbackslash An}, \texttt{\textbackslash C}, \texttt{\textbackslash Oo}, \texttt{\textbackslash Restriction}, \texttt{\textbackslash Sn}, \texttt{\textbackslash Z}, \texttt{\textbackslash Zz}, \texttt{\textbackslash rank}, \texttt{\textbackslash rquotient}, \texttt{\textbackslash sets}, \texttt{\textbackslash sgn}, \texttt{\textbackslash spn}, \texttt{\textbackslash std}, \texttt{\textbackslash triv}), with extra wrapper packages (bm, bbm, dsfont, xspace, cancel, mathtools). The delivered numbering is the unit's own. Assets: no retained span contains a figure environment, a graphics-inclusion command, a PDF-page inclusion, a table or a TikZ picture; no image file is copied into this package, the package declares no asset dependency and no image file is redistributed with it. Licence: version arXiv:2312.11092v3 of this article is distributed under the Creative Commons Attribution 4.0 International licence, as recorded in the arXiv licence metadata for this exact version and shown on the abstract page https://arxiv.org/abs/2312.11092v3. A full rights notice is delivered at licenses/arxiv-2312.11092-CC-BY-4.0.txt. Characters: every retained excerpt is pure ASCII, so the delivered engine, XeLaTeX with its default Latin Modern text font, reports no missing character.
\begin{lem}
\label{lem orthogonal idempotents for unextended set}
\begin{enumerate}
\item[(a)] 
Let $Y\owns y_1$ be a transitive $\Gamma$-set for a finite abelian group $\Gamma$, such that the permutation representation $\rho_{\mathrm{perm}}$ of $\Gamma$ decomposes into 
characters as $\C[Y]=\bigoplus_\rho\rho$. Then in $K_\Gamma(Y\times Y)$, the classes
%
\[
t_\rho:=\frac{1}{\# Y}\sum_{\gamma\in\Gamma}\rho(\gamma^{-1})[\Oo_{\
\Gamma\cdot (y_1,\gamma y_1)}]
\]
%
give a system of orthogonal idempotents in $K_\Gamma(Y\times Y)$
summing to the identity endomorphism $\Oo_{\Gamma\cdot(y_1,y_1)}$ of $K_\Gamma(Y)$.
%
\item[(b)]
For all $\rho$, $\rank(t_\rho, \C[Y])=1$.
\end{enumerate}
\end{lem}

\begin{equation}
\label{eqn trho mu transforms according to rho}
(t_{\rho}\mu)(\gamma x)=\rho(\gamma)(t_{\rho}\mu)(x).
\end{equation}

\begin{lem}
\label{lem trho act with rank 1 or 0}
Suppose that $Y$ is a transitive $\Zz$-set, where $\Gamma=\pi_0(\Zz)$ is finite abelian or is $\Sn_3$. Let 
$E_{s,\rho_1}$ be a simple $K_\Zz(Y\times Y)$-module such that $t_{\rho}E_{s,\rho_1}\neq 0$. Then $t_\rho$ acts 
as a rank 1 idempotent on $E_{s,\rho_1}$ if $\Gamma$ is abelian, and also if $\Gamma=\Sn_3$ and $\#Y\neq 3, 6$. If $\# Y=3,6$, we can find still find a family of orthogonal rank $1$ idempotents summing to the identity.
\end{lem}

\begin{proof}
First assume that $\Gamma$ is abelian.

As $s$ acts via $\Gamma$, either $Y^s=\emptyset$ or $Y^s=Y$. By assumption, 
we are not in the first case.

Now, the action of the centralizer of $s$ factors through 
%
\[
\pi_0(Z_{\Zz}(s)/Z(\Zz))\to\Gamma;
\]
%
let $\Gamma_1\subset\Gamma$ denote the image.

For any $\mu\colon Y\to\C$ transforming under $\Gamma_1$ by $\rho_1$, we have
% 
\begin{align*}
(t_{\rho}\mu)(x)
&=
\frac{1}{\# Y}
\sum_{\gamma\in\Gamma}\rho(\gamma^{-1})\left([\Oo_{\Gamma(y_1,\gamma y_1)}]\mu\right)(x)
\\
&=
\frac{1}{\# Y}
\sum_{\dot{\gamma}\in\Gamma_1\rquotient\Gamma}\sum_{\gamma_1\in\Gamma_1}
\rho(\dot{\gamma}^{-1}\gamma_1^{-1})\mu(\gamma_1\dot{\gamma} x)
\\
&=
\frac{1}{\# Y}
\sum_{\dot{\gamma}\in\Gamma_1\rquotient\Gamma}\sum_{\gamma_1\in\Gamma_1}
\rho(\dot{\gamma}^{-1})\left(\rho(\gamma_1^{-1})\rho_1(\gamma_1)\right)\mu(\dot{\gamma} x)
\\
&=
\begin{cases}
\frac{\#\Gamma'}{\# Y}\sum_{\dot{\gamma}\in\Gamma_1\rquotient\Gamma}\rho(\dot{\gamma}^{-1})\mu(\dot{\gamma}x)& \text{if}~\rho~\text{restricts to}~\rho_1
\\
0&\text{otherwise}
\end{cases}.
\end{align*}
%
%&=
%\frac{1}{\#\Gamma}
%\sum_{\dot{\gamma}\in\Gamma'\rquotient\Gamma}\rho(\dot{\gamma}^{-1})\mu(\dot{\gamma}x)\sum_{\gamma_1\in\Gamma'}
%\rho(\dot{\gamma}^{-1})\rho'(\gamma_1)
%\\
(Note that in the first case, the function $\dot{\gamma}\mapsto\rho(\dot{\gamma}^{-1})\mu(\dot{\gamma}x)$ is well-defined.)
Moreover, if $\rho$ does restrict to $\rho_1$, then as in \eqref{eqn trho mu transforms according to rho}, we have that 
$t_\rho\mu$ transforms under all of $\Gamma$ via $\rho$.
%
Thus
%
\[
E_{s,\rho_1}=\spn{\sets{\mu_{\rho}}{\rho\Restriction_{\Gamma_1}=\rho_1}},
\]
%
and
%
\[
\sets{t_\rho}{t_\rho E_{s,\rho_1}\neq 0}=\sets{t_\rho}{\rho\Restriction_{\Gamma_1}=\rho_1}.
\]
%
In particular, $\mathrm{rank}(t_\rho, E_{s,\rho_1})=1$ if it is nonzero, by 
Lemma \ref{lem orthogonal idempotents for unextended set} (b).


Now assume that $\Gamma=\Sn_3$.
First suppose $Y=3$. Then $Y^s$ is a point or $Y^s=Y$. 
If $Y^s$ is a point then $t_d$ itself has rank 1. If $Y^s=Y$, then $\C[Y]$ can decompose under $\pi_0(Z_{G^\vee}(u,s))\to\Sn_3$ as either $\triv\oplus\std$, the regular representation of $\An_3$, $\triv\oplus\triv\oplus\sgn$, or $\triv^{\oplus 3}$. If $\rho_1=\std$ or in the $\An_3$ case, then $t_d$ itself is again rank $1$, $\rho_1=\sgn$. If $\mathrm{rank}\pi(t_d)=2$, then $Y^s=Y$, $\pi_0(Z(u,s))=\Z/2\Z$, and $Y$ is a union of a point and a size two orbit. Taking the basis of functions $Y\to\triv$ given by the constant function and the indicator function of the size two orbit,  the action of the idempotents constructed above is given by the matrices
%
\begin{equation}
\label{eqn idempotent S3 two triv case}
t_{d,\triv}=\begin{pmatrix}
1&\frac{1}{2}\\
0&0
\end{pmatrix},~~~
t_{d,\std}=\begin{pmatrix}
0&\frac{-2}{3}\\
0&1
\end{pmatrix}.
\end{equation}
%
Hence these orthogonal idempotents each have rank 1 and give the required decomposition.

If $\mathrm{rank}(t_d, E_{s,\rho_1})=3$, then $\C[Y]=\triv^{\oplus 3}$. In this case, adding the indicator function of a point to the basis used in \eqref{eqn idempotent S3 two triv case} gives a basis in which
%
\[
\pi(t_{d,\std}-\frac{9}{2}\pi(t_{d,\std}t_{d,\triv})\colon \chi_Y\mapsto 0, \chi_{y_1,y_2}\mapsto \chi_{1,2}-\frac{2}{3}\chi_1, \chi_1\mapsto 0
\]
and
\[
\frac{9}{2}t_{\rho}t_{\triv}\colon \chi_1\mapsto \chi_1, \chi_Y,\chi_{1,2}\mapsto 0
\]
and $t_{d,\triv}$, $t_{d,\std}-\frac{9}{2}\pi(t_{d,\std}t_{d,\triv}$, and $\frac{9}{2}t_{\rho}t_{\triv}$ give three rank 1 orthogonal idempotents at $\chi_0$, hence at all $\chi$. As each has rank 1, and they sum to the identity, they give a decomposition into lines for all $\chi$.

Now suppose that $\# Y_d=6$. In this case $Y_d^s=\emptyset$ or $Y_d^s=Y$. Then we can have $\mathrm{rank}(t_d, E_{s,\rho_1})>2$ in the following cases: 

If the image of $\pi_0(Z_{G^\vee}(u,s))=\Sn_3$, and $\rho_1=\std$, then the two Young symmetrizers given above give the decomposition into lines.

If the image is $\An_3$, then $Y^s$ decomposes as two copies of the regular representation of $\An_3$ and $t_d$ has rank $2$.

If $\rho_1=\triv_{\An_3}$, then the orthogonal idempotents $t_{\triv_{\Sn_3}}$ and $t_{\sgn}$ both have rank $1$, as required. If $\rho_1$ is nontrivial, then as $\std|_{\An_3}$ decomposes with multiplicity one, the orthogonal idempotents $t_{\std}$, $t_{\std'}$ each have rank $1$.
%Indeed, if $f\star\mu=\mu$, then $f\star p\circ\mu=p\circ\mu$ for any linear map $p$.

If the image of $\pi_0(Z(u,s))$ is $\Z/2\Z$, then $\C[Y]=\triv^{\oplus 3}\oplus\sgn^{\oplus 3}$ and $\mathrm{rank}(
t_d, E_{s,\rho_1})=3$.
In either case each of $t_{\triv}$, $t_{\std}$, $t_{\std'}$, respectively $t_{\sgn}$, $t_{\std}$, and $t_{\std}$ have rank $1$, as $\std$ restricts to $\triv\oplus\sgn$.

If the image of $\pi_0(Z(u,s,))$ is trivial, then $\mathrm{rank}\pi(t_d)=6$ and elements of $t_dJ_ut_d$ act as $6\times 6$-matrices. In this case $t_{\std}$ and $t_{\std'}$ can be diagonalized by integer matrices, and we obtain again a full family of idempotents.
%
\end{proof}

\end{document}
